% Source: https://artofproblemsolving.com/community/c3h1316658_intro_to_alg_confusion
% Stable identity: aops-topic-1316658-q01
Multiplying by $x^2-5x+6=(x-2)(x-3)$, the required identity becomes
\[2x=A(x-2)+B(x-3)\qquad\text{for all }x\ne2,3 .\]
\textbf{Why we may substitute $x=2$ and $x=3$.} Both sides are polynomials of degree at most $1$. Their difference $p(x)=2x-A(x-2)-B(x-3)$ is a polynomial that is $0$ at every $x\ne2,3$, that is, at infinitely many points. A nonzero polynomial of degree at most $1$ has at most one root, so $p$ is the zero polynomial, and the identity holds for \emph{every} $x$, including $2$ and $3$. (Equivalently: both sides are continuous, so equality near $x=3$ forces equality at $x=3$.) The substitution is made in the cleared equation, which is a different equation from the original one, so no division by zero occurs.

Now $x=3$ gives $6=A\cdot1$, so $A=6$, and $x=2$ gives $4=B\cdot(-1)$, so $B=-4$. Hence
\[\boxed{A=6,\ B=-4},\qquad \dfrac{2x}{x^2-5x+6}=\dfrac{6}{x-3}-\dfrac{4}{x-2}.\]
Check with coefficients: $A+B=2$ and $-2A-3B=-12+12=0$, matching $2x+0$.
